
Les relations à prouver se vérifiant ponctuellement sur @@M0@@, on est ramené au cas @@M1@@ local, avec @@M2@@. Les @@M3@@ sont alors libres sur @@M4@@, et les relations @@M5@@ pour @@M6@@ équivalent aux relations @@M7@@ pour @@M8@@. Or on peut calculer les @@M9@@ en utilisant la résolution projective @@M10@@ de @@M11@@; ce sont les groupes d'homologie du complexe @@M12@@. Comme on suppose @@M13@@ régulier, ce dernier est acyclique en degrés @@M14@@, d'où les relations @@M15@@, pour @@M16@@.

On en déduit les relations @@M17@@ pour @@M18@@, en procédant par récurrence sur @@M19@@ grâce à la suite exacte des Tor associée à la suite exacte:
@@M20@@
Comme pour tout @@M21@@ on a @@M22@@, on tire de cette suite exacte et de @@M23@@ les relations @@M24@@. De la suite exacte:
@@M25@@
et de @@M26@@, on déduit les relations @@M27@@. Enfin, de la suite exacte des Tor associée à la suite exacte:
@@M28@@
on déduit les relations @@M29@@ pour @@M30@@, par récurrence sur @@M31@@, le cas @@M32@@ étant évident.

b) On va montrer maintenant que les @@M33@@ sont projectifs sur @@M34@@. La question étant locale sur @@M35@@, on peut supposer @@M36@@ libre, soit @@M37@@, de sorte que @@M38@@ est défini par une suite @@M39@@ d'éléments de @@M40@@. Posons @@M41@@ et soit @@M42@@ l'homomorphisme qui applique
